% The Kauffman parity/primality layer consumed by the compositional route.
%
% This section separates Kauffman's published theorem from the additional
% factor-to-ambient lifting assertion made in the Goertzel/Fable assembly.
% Every result carries a local status marker.

\section{The Kauffman trail spine and its exact composition seam}
\label{sec:kauffman-trail-spine}

The final step of the source route is expressed in the language of
Spencer--Brown formations and Kauffman trails.  Two logically different
claims have too often been bundled under the name "the Kauffman spine".
The first is Kauffman's published parity/primality theorem.  The second is a
new composition claim: if a factored trail has smaller colourable factors,
then completion sequences for those factors may be executed in the original
trail and glued.  The first has a published proof.  The second depends on the
precise support convention for the source's between-region moves and must be
proved for the framed trail object used by this development.

Our primary published reference for the first part is L.~H.~Kauffman,
"Reformulating the map color theorem", \emph{Discrete Mathematics} 302
(2005), 145--172, especially Sections~3--4.  This is the expanded and
clarified version of "On the map theorem", \emph{Discrete Mathematics}
229 (2001), 171--184.  The statements below are also the ones reproduced in
Section~3 of Goertzel's v23 manuscript.  The present section makes explicit
where the v24 assembly asks for more.

\subsection{Coloured trails and completion}

\begin{definition}[coloured trail]
\label{def:coloured-trail}
A \emph{coloured trail} consists of a plane subcubic graph $G(T)$, two
distinguished degree-two vertices $u,v$, a proper $3$-edge-colouring
$C$ of $G(T)$, and a designated missing edge $e=uv$.  The extended
graph $G^{*}(T)$ is obtained by inserting $e$.  In formation language,
the two distinguished boundary curves are the containers and the remaining
formation lies in the annular between-region.
\end{definition}

At a degree-two vertex exactly two colours occur, so there is a unique
missing colour.  Write $\mu_C(u)$ and $\mu_C(v)$ for those colours.

\begin{lemma}[one-edge completion criterion]
\label{lem:trail-one-edge-completion}
\Sverified{Mettapedia.GraphTheory.FourColor.GoertzelV24DeletedEdgeTrail.exists_taitColoring_of_framedTangleCompletable}
The colouring $C$ extends across $e$ precisely when
$\mu_C(u)=\mu_C(v)$.  More generally, if an allowed sequence of
between-region Kempe moves takes $C$ to a colouring $C'$ with equal
missing colours, then $G^{*}(T)$ has a Tait colouring.
\end{lemma}

\begin{proof}
The new edge must receive a colour absent at both endpoints, which proves
necessity.  Conversely, colour it by the common missing colour.  Properness
is unchanged away from $u,v$, and at each endpoint the three incident
colours are then distinct.  Kempe moves preserve properness, so the same
argument applies after an allowed move sequence.
\end{proof}

\begin{definition}[source completability]
\label{def:source-trail-completability}
The coloured trail $(T,C)$ is \emph{completable} if a finite sequence of
the source's allowed between-region moves takes $C$ to a colouring $C'$
with $\mu_{C'}(u)=\mu_{C'}(v)$.
\end{definition}

We must distinguish two uses of the word ``coloured''.  The edge-deleted
graph $G(T)$ already carries the displayed colouring $C$ by definition.  The
trail is \emph{extension-colourable} only when $G^{*}(T)$ admits a Tait
colouring.  The working manuscript's phrase ``every colorable trail is
completable'' means
\[
  G^{*}(T)\text{ is Tait-colourable}
  \quad\Longrightarrow\quad
  (T,C)\text{ is completable for the specified starting colouring }C.
\]
The converse follows from Lemma~\ref{lem:trail-one-edge-completion}.  This
terminology is load-bearing: deleting an edge from an uncolourable cubic
graph produces a \emph{coloured} $G(T)$ but not an extension-colourable
trail.  Consequently the Trail Completability Theorem cannot be applied
directly to the deleted-edge trail of a least counterexample; minimality is
first spent on its strict factors, whose extended graphs are smaller and
therefore extension-colourable.

The word "allowed" is load-bearing.  Kauffman's parity theorem concerns
switches on closed two-colour circuits.  The corrected framed trail calculus
also admits the defect paths forced by the two degree-two vertices.  Those
path moves are precisely outside the scope of closed-formation parity.  The
development has already proved that its framed move relation is the union of
the Kauffman circuit moves and these defect-path moves, subject to the
source/graph boundary alignment.
\Sverified{Mettapedia.GraphTheory.FourColor.GoertzelV24BetweenRegionDefectPaths.framedTangleLegalKempeStep_iff_kauffman_or_defectPath}

\subsection{The published parity theorem}

For a proper Tait colouring $C$ of a cubic graph $G$, let
$N_{ab}(C)$ be the number of components of the two-colour subgraph on
colours $a,b$, and define
\[
  \Delta(C)=N_{rb}(C)+N_{rp}(C)+N_{bp}(C),
  \qquad \pi(C)=\Delta(C)\pmod 2.
\]
Each two-colour subgraph is $2$-regular, so its components are circuits.

\begin{theorem}[Spencer--Brown--Kauffman parity lemma]
\label{thm:kauffman-parity}
\Sprose
If $G$ is a plane cubic graph and $C'$ is obtained from $C$ by
switching two colours on one two-colour circuit, then
$\pi(C')=\pi(C)$.
\end{theorem}

\begin{proof}
It is enough to work in the connected component containing the operated
circuit.  Such a component has no bridge: each colour class is a perfect
matching, whereas deleting a bridge in a cubic graph leaves an odd number
of vertices on either side, forcing every perfect matching to contain that
bridge.  Three disjoint colour classes cannot all contain it.  In particular
the two sides of every edge are distinct, and the face-cell accounting below
has no repeated bridge incidence.

Identify the three edge colours with the three nonzero elements
$a,b,c$ of the Klein group
$\mathbb K=\mathbb F_2^2$, so $a+b=c$.  We first integrate the edge
colouring on the dual map.  Fix one face $f_0$ and put $q(f_0)=0$.  For any
other face $f$, choose a dual walk from $f_0$ to $f$ and define $q(f)$ to be
the sum in $\mathbb K$ of the colours of the crossed primal edges.

This is independent of the chosen walk.  Indeed, over $\mathbb F_2$ the
cycle space of a spherical dual map is generated by its facial boundaries,
and a dual face is the star of a primal vertex.  The sum around such a face
is $a+b+c=0$, because the three edges incident with a cubic vertex have the
three distinct colours.  Thus every closed dual walk has colour sum zero.
Across a primal edge $e$ with incident faces $f,f'$ we consequently have
\[
                   q(f')-q(f)=C(e).
\]
In particular $q$ is a proper four-colouring of the faces.  The three faces
incident with any primal vertex have three distinct potentials.

For $x\in\{a,b,c\}$ let $X_x$ be the union of the closed faces whose
potentials lie in the dyad $\{0,x\}$.  Locally at a cubic vertex, one or two
of the three incident faces belong to $X_x$; hence $X_x$ is a compact
subsurface of the sphere with boundary.  If $\{x,y,z\}=\{a,b,c\}$, an edge
lies on $\partial X_x$ exactly when its two face potentials lie on opposite
sides of
\[
                 \{0,x\}\mid\{y,z\},
\]
which is exactly when its edge colour is $y$ or $z$.  Therefore
$\partial X_x$ is the $(y,z)$ two-colour subgraph and
\[
                 b(X_x)=N_{yz}(C).
\]
Every component of a subsurface of the sphere has genus zero, so
$\chi(X_x)=2k(X_x)-b(X_x)$.  Reducing modulo two gives
\[
 N_{bc}(C)\equiv\chi(X_a),\qquad
 N_{ac}(C)\equiv\chi(X_b),\qquad
 N_{ab}(C)\equiv\chi(X_c).                     \tag{1}
\]

We next add the three Euler characteristics cell by cell.  Every primal
vertex belongs to all three $X_x$, and therefore contributes $3\equiv1$.
A face of potential $0$ belongs to all three, while a face of nonzero
potential $x$ belongs only to $X_x$; hence the total face contribution is
$F(G)$ modulo two.  Finally, an edge incident with a zero-potential face
belongs to all three $X_x$, whereas an edge whose two incident face
potentials are distinct nonzero elements belongs to exactly two.  If
$d_0(C)$ denotes the sum of the boundary lengths of the zero-potential
faces, no edge is counted twice in $d_0(C)$ because $q$ is proper.  Thus
(1) yields the concrete parity formula
\[
                  \Delta(C)\equiv
                  V(G)+F(G)+d_0(C)\pmod2.       \tag{2}
\]

Now let $K$ be the $(a,b)$-circuit on which the switch is performed.
Combinatorial separation for spherical maps says that the edge set of the
cycle $K$ is the boundary of a set $S$ of faces: an edge belongs to $K$ if
and only if exactly one of its incident faces lies in $S$.  This is the
dual form of the fact that facial boundaries span the cycle space; it is
precisely Corollary~\ref{cor:separation}, so no unstated Jordan-curve step is
being used here.

Define
\[
 q'(f)=\begin{cases}
        q(f)+c,&f\in S,\\
        q(f),&f\notin S.
       \end{cases}
\]
Edges off $K$ have either both incident faces in $S$ or neither, so their
differences are unchanged.  Along $K$, adding $c$ at exactly one end sends
$a$ to $b$ and $b$ to $a$.  Hence $q'$ is the face potential of the
switched colouring $C'$.

Only the faces in $S$ can change their contribution to $d_0$.  Inside $S$,
translation by $c$ exchanges potentials $0$ and $c$ and exchanges $a$ and
$b$.  Consequently
\[
 d_0(C')-d_0(C)\equiv
 \sum_{\substack{f\in S\\q(f)\in\{0,c\}}}|\partial f|\pmod2.       \tag{3}
\]
In the sum on the right, an edge joining two selected faces is counted
twice.  The edges counted once form the boundary of the selected face set
$\{f\in S:q(f)\in\{0,c\}\}$.  Every such edge separates the dyad
$\{0,c\}$ from $\{a,b\}$, so it has colour $a$ or $b$.  At each cubic
vertex this boundary has degree zero or two; it is therefore a disjoint
union of alternating $(a,b)$ circuits.  Every such circuit has even length.
The right-hand side of (3) is thus zero, and
$d_0(C')\equiv d_0(C)\pmod2$.

The map itself has the same numbers of vertices and faces before and after
the switch.  Formula (2) now gives
$\Delta(C')\equiv\Delta(C)\pmod2$, as required.  This is the
Tutte-style form of the Spencer--Brown--Kauffman parity argument: planarity
enters exactly through face potentials and spherical cycle separation.
\end{proof}

Planarity is essential.  Kauffman's one-edge-deleted Petersen formation
changes the curve count from five to four under a simple operation.  A
nonplanar counterexample is therefore a regression fixture for any eventual
formal statement: a theorem that omits the spherical-map hypothesis is
false.

\subsection{Prime trails and Kauffman's equivalence}

In Kauffman's terminology, a same-context trail is \emph{factored} if,
after the two contextual curves are removed by idemposition, the remaining
trail structure has more than one component.  In the different-context
case it is factored if an extra formation curve misses both endpoints of the
empty edge.  A trail is \emph{prime} if no recolouring in its orbit has a
factored representative.

\begin{lemma}[Kauffman's curve-count dichotomy]
\label{lem:kauffman-curve-count}
\Sprose
Let $T$ be a nonempty prime uncolourable planar trail.
\begin{enumerate}
  \item In an unfactored representative whose contextual curves have the
        same colour, the total two-colour circuit count is five.
  \item In an unfactored representative whose contextual curves have
        different colours, the total two-colour circuit count is four.
\end{enumerate}
\end{lemma}

\begin{proof}
In the same-colour case, primeness leaves one trail curve between the two
contextual curves.  Each endpoint of the empty edge lies on an alternating
circuit.  These two circuits are distinct, since a single alternating
circuit through both endpoints is a two-colour completion path.  Any
additional circuit would remain as an extra component after removal of the
contexts and would factor the trail.  The count is therefore two context
curves, one trail curve, and two alternating circuits: $2+1+2=5$.

In the different-colour case, unfactoredness says that every curve meets an
endpoint of the empty edge.  At the pure context endpoint there is one
context curve and one alternator.  At the endpoint lying on the superposed
context there is one trail-colour curve and the second context curve.  Any
further curve misses both endpoints and is a factor.  Hence the count is
two context curves, one trail-colour curve, and one alternator:
$2+1+1=4$.  This is the incidence count in Kauffman's Section~4; its two
endpoint cases are definitions of the two context types, not an exhaustive
graph census.
\end{proof}

\begin{theorem}[Kauffman's primality equivalence]
\label{thm:kauffman-primality-equivalence}
\Sprose
For plane cubic graphs, the Four-Colour Theorem is equivalent to the
Primality Principle:
\[
  \text{every minimal nonempty uncolourable planar trail is prime}.
\]
\end{theorem}

\begin{proof}
Assume the Primality Principle and choose a minimal nonempty uncolourable
trail.  It is prime.  Begin with same-colour contexts, so
Lemma~\ref{lem:kauffman-curve-count} gives the odd count five.  A simple
operation on one context changes to the different-colour situation while
Theorem~\ref{thm:kauffman-parity} preserves parity.  An unfactored
different-colour prime trail would have the even count four.  The operated
formation is therefore factored, contradicting primeness.  Thus no such
trail exists.  Deleting an edge of a least uncolourable bridgeless cubic
plane graph produces a coloured nonempty trail whose extension is the
original graph, a contradiction.

Conversely, if the Four-Colour Theorem holds then no nonempty uncolourable
planar trail exists, so the Primality Principle holds vacuously.  This is
Kauffman's published equivalence argument.
\end{proof}

The finite "Closed Splitting" audits in the v24 working manuscript support
the parity behaviour on their listed fullerene instances.  They are not a
replacement for Theorem~\ref{thm:kauffman-parity}, nor do finite instances
prove Theorem~\ref{thm:kauffman-primality-equivalence}.

\subsection{The extra theorem required by the compositional route}

The source now proposes to prove its Trail Completability theorem by strong
induction.  In a factored minimal obstruction, each factor has a smaller
extended graph and is therefore colourable.  Applying Trail Completability
to each factor produces factor-local completion sequences.  One further
theorem is needed before those sequences may be composed in the original
formation.

\begin{lemma}[finite product reachability]
\label{lem:finite-product-reachability}
\Sprose
Let $I$ be finite.  For each $i\in I$, let $R_i$ be a relation on a state
space $X_i$, let $x_i\in X_i$, and suppose
$x_i\,R_i^{*}\,y_i$.  On $\prod_i X_i$, let $R_{\Pi}$ be the asynchronous
product relation: one step changes exactly one coordinate by an $R_i$-step
and fixes every other coordinate.  Then
\[
       (x_i)_{i\in I}\;R_{\Pi}^{*}\;(y_i)_{i\in I}.
\]
Consequently, if $y_i$ belongs to a local accepting set $A_i$ for every
$i$, the initial product state reaches $\prod_i A_i$.
\end{lemma}

\begin{proof}
Choose an order on the finite set $I$.  Lift the first coordinate's finite
$R_i$-path to product states by holding all other coordinates fixed, then
lift the second coordinate's path, and so on.  Concatenating these lifted
paths gives the required $R_{\Pi}^{*}$ path.  Previously completed
coordinates are never changed by a later lifted step.  The accepting-set
claim is immediate from the endpoint tuple.
\end{proof}

\begin{proposition}[source factorization has product semantics]
\label{prop:factor-product-semantics}
\Sopen
Every reachable Kauffman factorization of a source framed trail determines
factor state spaces $X_i$, factor move relations $R_i$, and a boundary
compatibility subset $K\subseteq\prod_i X_i$ such that:
\begin{enumerate}
  \item restriction maps every ambient reachable colouring to a point of
        $K$;
  \item every factor move lifts to the asynchronous product relation while
        preserving $K$, and every such lifted step is a legal ambient move;
  \item the product of the factor-completion accepting sets, intersected
        with $K$, maps to an ambient-completable colouring; and
  \item every factor is strict in the induction measure.
\end{enumerate}
\end{proposition}

This is the precise missing bridge between the topological word
"factored" and the algebraic word "independent".  Once it is proved,
Lemma~\ref{lem:finite-product-reachability} performs all scheduling and
commutation; no informal appeal to simultaneous moves remains.

\begin{proposition}[factor completion lifts and glues]
\label{prop:factor-completion-lifts}
\Sconditional{Proposition~\ref{prop:factor-product-semantics}}
Let a reachable representative of a coloured source trail $T$ be factored
into strict factors $T_1,\ldots,T_k$.  Suppose that, for every $i$, the
specific colouring induced on $T_i$ is completable by the source's framed
between-region moves.  Then there is a sequence of legal moves of $T$ whose
restriction realizes the required factor completions and after which $T$
itself is completable.
\end{proposition}

\begin{proof}
Apply Proposition~\ref{prop:factor-product-semantics}.  The factor
completion paths are paths in the relations $R_i$.  Lemma
\ref{lem:finite-product-reachability} combines them into an asynchronous
product path that stays in $K$.  Item~2 lifts that path to legal ambient
moves, and item~3 turns its accepting endpoint into an ambient completion.
\end{proof}

This is not a formal nicety.  A two-colour component that is closed in an
isolated factor can continue through a contextual curve when the factor is
placed back in the ambient formation.  Switching only its factor portion is
then not an ambient Kempe move.  Moreover, distinct factors share contextual
curves, so arbitrary completion sequences need not have disjoint supports.
The proposition must therefore be proved by one of the following exact
mechanisms, not by the sentence "the factors are disjoint":
\begin{enumerate}
  \item strengthen factor completability to sequences fixing the complete
        factor interface, then use disjoint-support commutation;
  \item prove that every factor move has an ambient component lift whose
        action on the other factors cancels or preserves their completion
        states; or
  \item replace sequential completion by a simultaneous product-state
        reachability theorem for all factors and their shared contexts.
\end{enumerate}

The current framed-trail API supplies frozen-interface vocabulary and exact
source/graph move alignment, but no theorem presently establishes any of
these three alternatives for a Kauffman factorization.  Proposition
\ref{prop:factor-product-semantics} is consequently ordinary mathematics
still to be done.  It must not be delegated to a finite census unless the
source first proves a uniform bound on the factor interface.

\begin{theorem}[source-faithful Kauffman assembly]
\label{thm:kauffman-source-assembly}
\Sconditional{Trail Completability and Proposition~\ref{prop:factor-completion-lifts}}
Assume the source's Trail Completability theorem for every strictly smaller
colourable factor, and assume Proposition~\ref{prop:factor-completion-lifts}.
Then every minimal nonempty uncolourable planar trail is prime.  Consequently
every bridgeless cubic plane graph is Tait colourable.
\end{theorem}

\begin{proof}
Let $T$ be a minimal nonempty uncolourable planar trail.  If it is not
prime, some reachable representative factors as
$T_1,\ldots,T_k$.  Each extended factor is strictly smaller than the
extended graph of $T$, so minimality makes it colourable.  Trail
Completability supplies a completion of its induced colouring.  Proposition
\ref{prop:factor-completion-lifts} glues those completions and makes $T$
completable, contradicting uncolourability.  Thus the Primality Principle
holds, and Theorem~\ref{thm:kauffman-primality-equivalence} gives the claimed
Tait colouring statement.
\end{proof}

Thus the Kauffman layer is not one undifferentiated open box.  Its parity
theorem, curve-count dichotomy, and equivalence are established ordinary
mathematics.  The live source-specific seam is the ambient lifting and gluing
of factor completion sequences.  That is the theorem the Lean layer should
formalize only after its ordinary proof has been supplied.
